% ============================================================================
%  Section VII — Discussion and Limitations  (IEEE TIM, IEEEtran)
%  Passive voice / third person throughout. Interpretation of the Section VI
%  results, practical guidance, and limitations. Requires: none beyond base.
% ============================================================================

\section{Discussion and Limitations}
\label{sec:discussion}

\subsection{Discussion}

\subsubsection{Evaluation methodology}
The optimism gap of Section~\ref{subsec:gap} is the central methodological
finding: for low-rate smartphone-based machine diagnosis, accuracies obtained under
random splitting of overlapping windows overstate deployable performance by tens
of accuracy points. The effect mirrors that documented for industrial-grade
vibration data~\cite{wheat2024impact, hendriks2022towards} and the broader
machine-learning leakage literature~\cite{kapoor2023leakage}, but its magnitude
here is large because each operating condition is represented by a single
continuous recording. A practical consequence is that the near-perfect accuracies
frequently reported for this class of data should be interpreted with caution; the
recording-wise and cross-condition figures obtained here provide a more
deployment-honest reference against which future methods can be compared.

\subsubsection{Low sampling rate is not a barrier}
Decimation to $25$~Hz did not reduce accuracy and in several runs improved it, so
the low sampling rate of consumer MEMS sensors is not a barrier to coarse-fault
discrimination. This observation warrants care rather than celebration: the sweep
is non-monotonic---the intermediate $50$~Hz rate performed worst---and, because
decimation applies an anti-alias low-pass, the gain at $25$~Hz is entangled with
the removal of high-frequency content that the classifier would otherwise overfit
on the small training set, rather than with the sample rate as such. The
defensible conclusion is the weaker but still useful one that the discriminative
information for these faults resides in the low-frequency band identified in
Section~\ref{sec:charact}, so that aggressive decimation reduces telemetry and, by
extension, node energy without sacrificing accuracy. On this dataset the
resource-frugal configuration established in Section~\ref{subsec:tradeoff}
---$25$~Hz sampling, the three raw channels, and a $4$~s window---is a strong
operating point; confirming it as a general default would require the
multi-machine validation identified in Section~\ref{subsec:limits}.

\subsubsection{Model complexity and edge feasibility}
Under leakage-free evaluation the lower-capacity models generalised at least as
well as the Random Forest while occupying a kilobyte-scale footprint, so on-node
inference is practical on inexpensive microcontrollers. Combined with the link
budget, this yields per-link deployment guidance: Bluetooth Low Energy, Zigbee,
and NB-IoT nodes can stream raw or feature data to a gateway, whereas LoRaWAN
nodes, constrained by the duty cycle to the decision stream, must perform
classification locally. The alignment between the models that deploy most easily
and those that generalise best is convenient for practical systems.

\subsubsection{Sensor-level guidance}
The channel analysis of Section~\ref{sec:charact} carries a concrete
recommendation for smartphone-based acquisition: the raw acceleration channels
should be preferred over the gravity-compensated channels, because the on-device
gravity compensation---valuable for motion-tracking applications---attenuates
low-frequency content that is diagnostically useful for condition monitoring.

\subsubsection{Scope of applicability}
Taken together, the results delineate where low-cost, low-rate MEMS sensing is
appropriate. Coarse health states, such as a healthy machine, a broken rotor bar,
or a supply-voltage imbalance, are separable at $100$~Hz and below; the staging of
incipient bearing degradation, whose signatures lie in the high-frequency band, is
not, and remains the province of higher-rate, industrial-grade instrumentation.

\subsection{Limitations}
\label{subsec:limits}
Several limitations bound the scope of these findings. The dataset comprises a
single $1.1$~kW machine with laboratory-induced, discrete-severity faults, and the
two lubrication states were graded qualitatively; the reported accuracies are
therefore a deployment-honest baseline rather than a claim of field performance.
Crucially, each fault class is represented by a \emph{single physical specimen}
(one modified bearing, one drilled rotor, one imbalance configuration) recorded
across speeds and loads. The recording-wise protocol therefore removes
window-overlap and exact operating-point leakage, but not specimen identity: a
model may still key on the idiosyncratic signature of that particular specimen and
its mounting rather than on a fault signature that would transfer to a different
specimen of the same class. The evaluation is thus leakage-free only in this
restricted sense, and cross-specimen generalisation is not established here. The
small number of recordings ($36$) also produces wide fold-to-fold variability;
although five-seed repetition stabilises the central estimates and the reported
expanded uncertainty quantifies their reproducibility, it captures only resampling
variance and not the dominant specimen-, machine-, and mounting-to-mounting
components, so it should not be read as a full measurement-uncertainty budget. The
deployment quantities are likewise indicative rather than measured: the model
footprint is a serialised (software) size, and energy and latency are inferred
from data volume and window length rather than measured on hardware, with a
cycle-accurate microcontroller characterisation left to future work. Finally,
generalisation across machines of different ratings and manufacturers was not
assessed and constitutes the most important direction for future validation.
